\question[3]
Ein mit Vigenère verschlüsselter Geheimtext lautet:

\texttt{HBWVHBWVHBWV}

Man kennt weder den Klartext noch den Schlüssel. Erläutere die Grundidee 
des Kasiski-Tests und wende sie auf das obige Beispiel an, um eine 
wahrscheinliche Schlüssellänge zu bestimmen.

\begin{solutionbox}{4.5cm}

Der Kasiski-Test nutzt aus, dass sich Buchstabenfolgen im Klartext oft 
wiederholen (z.\,B. häufige Wörter oder Silben). Trifft eine solche 
Wiederholung im Klartext an zwei Stellen auf dieselbe Phase des 
periodisch wiederholten Schlüssels, entsteht an beiden Stellen 
dieselbe Geheimtextfolge. Der Abstand zwischen zwei gleichen 
Geheimtextfolgen ist dann ein Vielfaches der Schlüssellänge.

Im Beispiel wiederholt sich die Folge \texttt{HBWV} an den Positionen 
1, 5 und 9. Die Abstände zwischen den Wiederholungen betragen 4 und 8 
Zeichen. Der größte gemeinsame Teiler dieser Abstände ist 4. Die 
wahrscheinliche Schlüssellänge beträgt daher \texttt{4}.

\end{solutionbox}


